\section{Manual técnico} \label{sec:appendix_B_tech_manual}

Este apéndice describe la organización técnica del proyecto y recoge una guía
práctica para la extensión de la aplicación con nuevas operaciones de procesado y nuevos
cargadores de formatos de imagen.

\subsection{Objetivo del Manual Técnico}

El sistema se ha diseñado para que la incorporación de nuevas funcionalidades no
requiera modificar de forma transversal toda la aplicación. En particular, dos de
los puntos de extensión más importantes son las operaciones de procesado y los
\textit{loaders} de imagen. Ambos siguen una estructura común basada en clases
base, modelos compartidos y registros centralizados.

\noindent{El propósito de este manual técnico es servir como guía para:}

\begin{itemize}[noitemsep]
    \item Localizar los módulos principales en el código fuente.
    \item Comprender el flujo de integración de una operación de procesado.
    \item Comprender el flujo de integración de un nuevo formato de imagen.
    \item Mantener la separación entre núcleo de procesamiento, carga de datos, interfaz gráfica y \textit{benchmarking}.
    \item Facilitar que el proyecto pueda evolucionar sin comprometer su mantenibilidad.
\end{itemize}

Esta capacidad de extensión es una de las principales ventajas de la
arquitectura desarrollada. Una nueva operación puede añadirse implementando una
clase concreta y registrando su definición visual, mientras que un nuevo formato
de imagen puede incorporarse mediante un \textit{loader} específico sin alterar la
lógica general del controlador.

\newpage

\subsection{Estructura general del proyecto}

El código fuente se organiza en torno a varios módulos principales que reflejan las
diferentes responsabilidades del sistema. Esta división permite que cada parte evolucione 
de forma independiente y reduce la probabilidad de introducir dependencias innecesarias.

\begin{itemize}[noitemsep]
    \item \path{core/}: contiene las operaciones de procesamiento de imagen: filtrado, color, geometría, normalización, formato y procesado Bayer.
    \item \path{core/image_operation.py}: define la clase base \texttt{ImageOperation}, de la que heredan las operaciones concretas.
    \item \path{core/registry.py}: mantiene el registro de operaciones disponibles mediante el decorador \\ \texttt{@register\_operation}.
    \item \path{loaders/}: contiene las clases encargadas de la carga de imágenes desde distintos formatos.
    \item \path{loaders/base_loader.py}: define la interfaz común \texttt{ImageLoader}.
    \item \path{loaders/image_loader.py}: centraliza la selección del \textit{loader} adecuado a partir de la extensión del fichero.
    \item \path{models/}: contiene modelos compartidos como \texttt{Image}, \texttt{Metadata} y \texttt{OperationDefinition}.
    \item \path{controller.py}: coordina la carga de imágenes, la ejecución de operaciones y la actualización de la interfaz.
    \item \path{gui/}: implementa la interfaz gráfica y los controles asociados a las operaciones.
    \item \path{benchmark/}: contiene la infraestructura para ejecutar perfiles de operaciones y obtener métricas de rendimiento.
    \item \path{tests/}: contiene pruebas funcionales de operaciones, formatos, geometría, color y normalización.
\end{itemize}

La idea fundamental es que el controlador no necesita conocer los detalles
internos de cada operación ni de cada formato de fichero. Para las operaciones,
trabaja con nombres registrados en el registro global de operaciones. Para las imágenes,
delega la carga en \path{loaders/image_loader.py}, que selecciona el \textit{loader}
adecuado en función de la extensión.

\newpage

\subsection{Cómo añadir una nueva operación de procesado}

Las operaciones de procesado se implementan como clases que heredan de
\texttt{ImageOperation}. Esta clase base define el método público \texttt{\_\_call\_\_},
comprueba que la entrada y la salida son tensores de imagen válidos y mide el
tiempo de ejecución de la operación. La clase concreta sólo debe implementar el
método \texttt{apply}.

\noindent{El flujo recomendado para añadir una nueva operación es el siguiente:}

\begin{enumerate}[noitemsep]
    \item Crear un nuevo fichero en el subpaquete adecuado de \path{core/}
    \item Definir una clase que herede de \texttt{ImageOperation}
    \item Decorar la clase con \texttt{@register\_operation}
    \item Implementar el método \texttt{apply(self, x: Tensor) -> Tensor}
    \item Validar los parámetros en el constructor de la operación
    \item Añadir una entrada en \path{models/operation_definition.py} para que la operación aparezca en la interfaz gráfica
    \item Añadir pruebas en \path{tests/} cuando la operación tenga comportamiento no trivial
\end{enumerate}

El \autoref{lst:technical_manual_operation_example} muestra un ejemplo reducido
de operación de inversión de intensidad para imágenes representadas como tensores
reales en el rango \texttt{[0, 1]}.

\begin{lstlisting}[style=pythonstyle, caption={Ejemplo simplificado de nueva operación de procesado}, label={lst:technical_manual_operation_example}]
from torch import Tensor

from core import _tensor_utils as T_u
from core.image_operation import ImageOperation
from core.registry import register_operation


@register_operation
class InvertIntensity(ImageOperation):
    def apply(self, x: Tensor) -> Tensor:
        T_u.assert_real_valued_tensor(x)
        return 1.0 - x
\end{lstlisting}

El decorador \texttt{@register\_operation} introduce la clase en el registro global
de operaciones. De esta forma, el controlador puede instanciarla a partir de su
nombre sin depender directamente del módulo donde está implementada.

\newpage

Si la operación no tiene parámetros, basta con añadir una definición de tipo
\texttt{no\_param}, como se muestra en el
\autoref{lst:technical_manual_operation_definition_no_param}.

\begin{lstlisting}[style=pythonstyle, caption={Definición de operación sin parámetros}, label={lst:technical_manual_operation_definition_no_param}]
OperationDefinition(
    name="InvertIntensity",
    label="Invert Intensity",
    category="Filters",
    control_type="no_param",
)
\end{lstlisting}

Si la operación requiere parámetros numéricos, se recomienda definir una lista de
\texttt{FloatParamSpec}. Estos objetos describen el nombre interno del parámetro,
su etiqueta visible, el rango permitido, el paso de ajuste y el valor por defecto.
El \autoref{lst:technical_manual_operation_definition_params} muestra un ejemplo
con un parámetro \texttt{factor}.

\begin{lstlisting}[style=pythonstyle, caption={Definición de operación con parámetros}, label={lst:technical_manual_operation_definition_params}]
_FILTER_EXAMPLE_PARAMS = [
    FloatParamSpec(
        key="factor",
        label="Factor",
        minimum=0.0,
        maximum=2.0,
        step=0.01,
        default=1.0,
    ),
]

OperationDefinition(
    name="ExampleOperation",
    label="Example Operation",
    category="Filters",
    control_type="filter",
    params=_FILTER_EXAMPLE_PARAMS,
    default_params=_build_filter_params(_FILTER_EXAMPLE_PARAMS),
)
\end{lstlisting}

Es importante que el valor de \texttt{name} coincida exactamente con el nombre de
la clase registrada. Esta coincidencia permite que el perfil de operaciones, la
interfaz gráfica y el módulo de \textit{benchmarking} utilicen el mismo
identificador. Así, una operación añadida correctamente queda disponible en la
interfaz, puede exportarse dentro de un perfil y puede ser ejecutada por el
sistema de medición de rendimiento.

En caso de requerir un control gráfico específico, debe añadirse un nuevo tipo de
control en \path{gui/right_panel/widgets/} y conectarlo desde el panel
derecho. Sin embargo, muchas operaciones pueden reutilizar controles existentes,
especialmente las operaciones sin parámetros o aquellas basadas en parámetros
numéricos simples.

\newpage

\subsection{Cómo añadir un nuevo \textit{loader} de formato de imagen}

Los \textit{loaders} encapsulan la lógica de lectura de cada formato de imagen.
Todos heredan de la clase base \texttt{ImageLoader}. Esta interfaz obliga a declarar
las extensiones soportadas, los formatos asociados y el método \texttt{load}. El
resultado debe ser siempre una instancia del modelo \texttt{Image}, que contiene el
tensor de imagen, su copia original, el formato, el espacio de color y los metadatos.

\noindent{El flujo recomendado para añadir un nuevo formato es el siguiente:}

\begin{enumerate}[noitemsep]
    \item Crear una nueva clase en \path{loaders/} heredando de \texttt{ImageLoader}
    \item Declarar la propiedad \texttt{extensions} con las extensiones soportadas en minúsculas
    \item Declarar la propiedad \texttt{formats} con los valores correspondientes de \texttt{ImageFormat}
    \item Implementar \texttt{load(path, device)} devolviendo un objeto \texttt{Image} completo
    \item Registrar el nuevo \textit{loader} en la lista \texttt{\_LOADERS} de \texttt{loaders/image\_loader.py}
    \item Añadir el nuevo valor de \texttt{ImageFormat} si el formato no existe todavía
    \item Añadir pruebas de carga con una imagen mínima representativa del formato
\end{enumerate}

\newpage

El \autoref{lst:technical_manual_loader_example} muestra la estructura básica de
un \textit{loader}. El ejemplo es muy genérico, la función de lectura
concreta dependerá del formato y de la biblioteca utilizada.

\begin{lstlisting}[style=pythonstyle, caption={Estructura simplificada de un nuevo loader}, label={lst:technical_manual_loader_example}]
import torch
from torch import Tensor

from loaders.base_loader import ImageLoader
from models.enums.color_space import ColorSpace
from models.enums.image_formats import ImageFormat
from models.image import Image
from models.metadata import Metadata


class TiffLoader(ImageLoader):
    @property
    def extensions(self) -> list[str]:
        return ["tif", "tiff"]

    @property
    def formats(self) -> list[ImageFormat]:
        return [ImageFormat.TIFF]

    def load(self, path: str, device) -> Image:
        tensor: Tensor = read_tiff_as_tensor(path)
        tensor = tensor.to(dtype=torch.float32, device=device)

        if not tensor.is_contiguous():
            tensor = tensor.contiguous()

        return Image(
            tensor=tensor,
            original_tensor=tensor.clone(),
            operation_result_tensor=tensor.clone(),
            path=path,
            name=path.split("/")[-1],
            image_format=ImageFormat.TIFF,
            color_space=ColorSpace.RGB,
            metadata=Metadata(
                width=tensor.shape[2],
                height=tensor.shape[1],
                bit_depth=16,
                bayer_pattern=None,
            ),
        )
\end{lstlisting}

\newpage

Una vez implementado el \textit{loader}, debe añadirse a la lista de cargadores
disponibles. El \autoref{lst:technical_manual_loader_registry} muestra el punto
de integración.

\begin{lstlisting}[style=pythonstyle, caption={Registro de un nuevo loader}, label={lst:technical_manual_loader_registry}]
from loaders.tiff_loader import TiffLoader

_LOADERS: list[ImageLoader] = [
    RawLoader(),
    ArwLoader(),
    DngLoader(),
    JpgLoader(),
    PklGzLoader(),
    TiffLoader(),
]
\end{lstlisting}

Con esta modificación, el sistema incorpora automáticamente las extensiones del
nuevo \textit{loader} al mapa interno de extensiones. La interfaz gráfica también
se beneficia de esta integración, ya que el filtro de selección de archivos se
construye a partir de las extensiones soportadas por los \textit{loaders}
registrados.

\newpage

\subsection{Recomendaciones de mantenimiento}

Para conservar la extensibilidad del proyecto, las ampliaciones futuras deberían
seguir las mismas reglas de separación aplicadas durante el desarrollo del
sistema. Las recomendaciones principales son las siguientes:

\begin{itemize}[noitemsep]
    \item Mantener las operaciones de procesado dentro de \texttt{core/}, evitando dependencias directas con la interfaz gráfica.
    \item Validar parámetros en el constructor de cada operación para detectar configuraciones inválidas antes de ejecutar el procesamiento.
    \item Usar \texttt{@register\_operation} en toda operación nueva para mantener el mecanismo de descubrimiento centralizado.
    \item Mantener sincronizados el nombre de la clase, el campo \texttt{name} de \texttt{OperationDefinition} y los perfiles exportados.
    \item Reutilizar controles gráficos existentes siempre que sea posible antes de crear nuevos \textit{widgets}.
    \item Encapsular la lectura de formatos en \textit{loaders} independientes, sin introducir lógica específica de formato en el controlador.
    \item Devolver siempre objetos \texttt{Image} completos desde los \textit{loaders}, incluyendo \texttt{Metadata} y \texttt{ColorSpace}.
    \item Asegurar que los tensores generados son contiguos y tienen la forma esperada por el núcleo de procesamiento.
    \item Añadir pruebas unitarias o de integración para toda operación o formato nuevo que pueda afectar al flujo principal.
    \item Mantener actualizado el perfil de \textit{benchmarking} cuando una nueva operación deba evaluarse en términos de rendimiento.
\end{itemize}
